% !TEX program = latexmk
% !TEX encoding = UTF-8 Unicode

\documentclass[11pt,a4paper]{article}

% --- Page / typography ---
\usepackage[a4paper,margin=2.5cm]{geometry}
\usepackage[T1]{fontenc}
\usepackage[utf8]{inputenc}
\usepackage{lmodern}
\usepackage{microtype}

% --- Math ---
\usepackage{amsmath,amssymb,amsthm}
\usepackage{mathtools}

% --- Figures / tables ---
\usepackage{graphicx}
\usepackage{booktabs}
\usepackage{float}
\usepackage{subcaption}

% --- Code listings ---
\usepackage{listings}
\usepackage{xcolor}

\lstset{
  basicstyle=\ttfamily\small,
  backgroundcolor=\color{gray!10},
  frame=single,
  breaklines=true,
  columns=fullflexible,
  keywordstyle=\color{blue},
  commentstyle=\color{gray},
  stringstyle=\color{orange}
}

% --- Links ---
\usepackage[hidelinks]{hyperref}
\usepackage{url}

% --- Graphics path ---
\graphicspath{{img/}}

% --- Algorithm environments ---
\usepackage{algorithm}
\usepackage{algpseudocode}

% --- Theorem environments ---
\theoremstyle{plain}
\newtheorem{theorem}{Theorem}[section]
\newtheorem{lemma}[theorem]{Lemma}

\theoremstyle{definition}
\newtheorem{definition}[theorem]{Definition}

\theoremstyle{remark}
\newtheorem*{remark}{Remark}

% --- Document metadata ---
\title{PCG Arena: A Web Platform for Blind A/B Testing of Procedural Content Generators}
\author{Antoni Czapski}
\date{February 2026}

\begin{document}
\maketitle

\begin{abstract}
Evaluating procedural content generators (PCGs) for video games poses significant methodological challenges, as player enjoyment is inherently subjective and difficult to quantify through automated metrics alone. This report presents \textbf{PCG Arena}, a web-based platform designed for blind A/B testing of Mario level generators through human preference data. The platform enables players to play two procedurally generated levels side-by-side in their browser, vote for their preferred level, and optionally tag gameplay characteristics. Votes are aggregated using the Glicko-2 rating system, which provides uncertainty-aware rankings with confidence intervals.

Key technical contributions include: (1) a browser-based Mario game engine faithfully ported from Java to TypeScript, (2) an adaptive matchmaking algorithm (AGIS) that optimizes for rapid rating convergence, (3) a builder profile system enabling researchers to submit and compete with their own generators, and (4) comprehensive research analytics including death heatmaps, player trajectories, and confusion matrices.

The platform is deployed at \url{https://pcg-arena.com}, featuring Google OAuth and email/password authentication, Twilio SendGrid integration for email verification, and Docker-based deployment on Google Cloud Platform. This work provides the research community with an accessible tool for empirically comparing procedural content generators through human evaluation.
\end{abstract}

\tableofcontents
\newpage

%==============================================================================
\section{Introduction}
\label{sec:introduction}
%==============================================================================

Procedural Content Generation (PCG) has become a cornerstone of modern game development, enabling the automatic creation of game levels, narratives, textures, and other assets~\cite{summerville2018}. However, evaluating the quality of procedurally generated content remains a fundamental challenge in the field. While automated metrics such as playability, diversity, and structural features provide valuable insights, they often fail to capture the subjective experience of human players---the ultimate consumers of generated content.

The Mario AI Championship's Level Generation Track~\cite{marioai} established a precedent for human evaluation of PCG by having players experience and compare levels from different generators. However, such competitions typically rely on controlled laboratory settings with limited participants, making it difficult to gather statistically significant preference data at scale.

This project addresses these limitations by developing \textbf{PCG Arena}, a fully operational web platform that enables:

\begin{enumerate}
    \item \textbf{Browser-Based Gameplay}: Players can experience Mario levels directly in their browser without downloading any software, removing the primary barrier to participation.
    
    \item \textbf{Blind A/B Testing}: Each battle presents two levels from different generators without revealing their source, eliminating bias from generator reputation.
    
    \item \textbf{Uncertainty-Aware Rankings}: The Glicko-2 rating system provides not only point estimates of generator quality but also confidence intervals that reflect the reliability of those estimates.
    
    \item \textbf{Adaptive Matchmaking}: The AGIS (Adaptive Glicko-Informed Selection) algorithm prioritizes matchups that maximize information gain, ensuring efficient convergence of ratings.
    
    \item \textbf{Research Analytics}: Comprehensive data collection enables analysis of player preferences, level difficulty, and engagement patterns.
\end{enumerate}

The platform is live at \url{https://pcg-arena.com} and serves as both a practical tool for PCG researchers and a demonstration of modern web technologies for game-based research.

This report is structured as follows: Section~\ref{sec:background} provides background on PCG and rating systems. Section~\ref{sec:architecture} describes the system architecture. Sections~\ref{sec:engine}--\ref{sec:admin} detail the core components. Section~\ref{sec:deployment} covers infrastructure, and Section~\ref{sec:results} presents preliminary results.

%==============================================================================
\section{Background and Related Work}
\label{sec:background}
%==============================================================================

\subsection{Procedural Content Generation}
\label{subsec:pcg}

Procedural Content Generation refers to the algorithmic creation of game content with limited or indirect human input~\cite{summerville2018}. PCG techniques range from simple rule-based systems to sophisticated machine learning approaches, including:

\begin{itemize}
    \item \textbf{Search-Based PCG}: Using evolutionary algorithms to optimize levels against fitness functions measuring playability, difficulty, and aesthetics.
    \item \textbf{Grammar-Based PCG}: Defining formal grammars that specify valid level structures and generating content through grammar expansion.
    \item \textbf{Machine Learning PCG (PCGML)}: Training neural networks on existing human-designed content to learn implicit design patterns~\cite{summerville2018}.
\end{itemize}

A fundamental challenge in PCG research is evaluation: how do we determine whether one generator produces ``better'' content than another? This question motivates the development of PCG Arena.

\subsection{Mario AI Framework}
\label{subsec:mario-ai}

The Mario AI Framework~\cite{marioai} provides a standardized environment for AI research in platformer games. Originally developed for the Mario AI Competition, it offers:

\begin{itemize}
    \item A faithful recreation of Super Mario Bros. physics and mechanics
    \item A standardized ASCII level format for representing tile-based levels
    \item APIs for both playing agents and level generators
    \item Tools for level analysis and visualization
\end{itemize}

PCG Arena builds upon this framework, using its level format and adapting its physics engine for browser-based gameplay. The ASCII tilemap format represents levels as grids of characters, where each character denotes a specific tile type (ground, brick, enemy spawn, etc.).

\subsection{Human Evaluation in PCG}
\label{subsec:human-eval}

Human evaluation remains the gold standard for assessing game content quality~\cite{marioai}. Common approaches include:

\begin{itemize}
    \item \textbf{Absolute Rating}: Players rate individual levels on a numerical scale (e.g., 1--5 stars).
    \item \textbf{Pairwise Comparison}: Players compare two items and choose their preference.
    \item \textbf{Ranking}: Players order multiple items from best to worst.
\end{itemize}

Pairwise comparison, the method used by PCG Arena, offers several advantages: it reduces cognitive load compared to absolute rating, produces more consistent results across participants, and maps naturally to ranking algorithms like Elo and Glicko-2.

\subsection{Rating Systems for Pairwise Comparisons}
\label{subsec:rating-systems}

Rating systems for pairwise comparisons originated in chess and have been widely adopted for competitive games and recommendation systems. The Elo rating system, developed by Arpad Elo, models the probability of player A defeating player B based on their rating difference:

\begin{equation}
P(A \text{ wins}) = \frac{1}{1 + 10^{(R_B - R_A)/400}}
\end{equation}

The Glicko-2 system~\cite{glickman2012} extends Elo by adding two additional parameters:

\begin{itemize}
    \item \textbf{Rating Deviation (RD)}: Measures uncertainty in the rating estimate. New or inactive players have high RD; active players converge to low RD.
    \item \textbf{Volatility ($\sigma$)}: Captures the expected fluctuation in performance over time.
\end{itemize}

PCG Arena uses Glicko-2 because its uncertainty quantification directly addresses a key research need: knowing not just which generator is best, but how confident we can be in that ranking.

%==============================================================================
\section{System Architecture}
\label{sec:architecture}
%==============================================================================

\subsection{High-Level Overview}
\label{subsec:arch-overview}

PCG Arena follows a classic three-tier web architecture with clear separation of concerns:

\begin{enumerate}
    \item \textbf{Presentation Layer}: React/TypeScript frontend running in the user's browser
    \item \textbf{Application Layer}: Python/FastAPI backend handling business logic and API endpoints
    \item \textbf{Data Layer}: SQLite database persisting all platform state
\end{enumerate}

Figure~\ref{fig:architecture} illustrates the high-level system architecture. Users interact with the platform through a browser-based frontend that communicates with the backend via a RESTful API following the \texttt{arena/v0} protocol specification.

\begin{figure}[H]
    \centering
    \includegraphics[width=0.9\textwidth]{admin_overview.png}
    \caption{Admin overview showing system architecture and platform statistics.}
    \label{fig:architecture}
\end{figure}

\subsection{Backend Architecture}
\label{subsec:backend}

The backend is implemented in Python using the FastAPI framework, chosen for its:

\begin{itemize}
    \item Automatic OpenAPI documentation generation
    \item Native async/await support for concurrent request handling
    \item Built-in request validation via Pydantic models
    \item High performance comparable to Node.js and Go
\end{itemize}

Key backend modules include:

\begin{itemize}
    \item \texttt{main.py}: FastAPI application with route definitions
    \item \texttt{auth.py}: Authentication logic (Google OAuth, email/password)
    \item \texttt{glicko2.py}: Glicko-2 rating system implementation
    \item \texttt{matchmaking.py}: AGIS matchmaking algorithm
    \item \texttt{builders.py}: Generator submission and management
    \item \texttt{stats.py}: Research analytics and data export
\end{itemize}

The backend runs in a Docker container, ensuring consistent deployment across environments.

\subsection{Frontend Architecture}
\label{subsec:frontend}

The frontend is built with React 18 and TypeScript 5.6, using Vite as the build tool. Key architectural decisions include:

\begin{itemize}
    \item \textbf{Component-Based Design}: UI elements are reusable React components
    \item \textbf{Canvas-Based Rendering}: Game graphics use HTML5 Canvas for direct pixel manipulation
    \item \textbf{Type Safety}: TypeScript catches errors at compile time
    \item \textbf{Fast Development}: Vite provides sub-second hot module replacement
\end{itemize}

The frontend directory structure separates concerns:

\begin{lstlisting}
frontend/src/
  api/           # Backend communication
  engine/        # Game engine (ported from Java)
  components/    # React UI components
  pages/         # Full page components
  contexts/      # React state management
\end{lstlisting}

\subsection{Database Design}
\label{subsec:database}

PCG Arena uses SQLite for data persistence, chosen for:

\begin{itemize}
    \item Single-file storage simplifying backups and deployment
    \item ACID transactions ensuring data consistency
    \item Zero configuration and no separate database service
    \item Sufficient performance for the expected load
\end{itemize}

The database schema comprises 17 tables across 17 migrations:

\begin{table}[H]
\centering
\begin{tabular}{ll}
\toprule
\textbf{Table} & \textbf{Purpose} \\
\midrule
\texttt{generators} & PCG algorithm metadata with owner reference \\
\texttt{levels} & ASCII tilemap levels with content hash \\
\texttt{battles} & Comparison instances pairing two levels \\
\texttt{votes} & User outcomes with telemetry data \\
\texttt{ratings} & Glicko-2 ratings per generator \\
\texttt{rating\_events} & Audit log of rating changes \\
\texttt{users} & Authenticated user accounts \\
\texttt{user\_sessions} & Active login sessions \\
\texttt{play\_trajectories} & Player position and death data \\
\bottomrule
\end{tabular}
\caption{Key database tables in PCG Arena.}
\label{tab:db-tables}
\end{table}

%==============================================================================
\section{Game Engine Implementation}
\label{sec:engine}
%==============================================================================

\subsection{TypeScript Port from Java}
\label{subsec:ts-port}

A critical component of PCG Arena is the browser-based Mario game engine, which was ported from the Java-based Mario AI Framework to TypeScript. This port required careful attention to maintain gameplay fidelity while adapting to browser constraints.

The porting process involved:

\begin{enumerate}
    \item \textbf{Class Structure Translation}: Java classes were converted to TypeScript classes, preserving inheritance hierarchies and encapsulation.
    \item \textbf{Physics Constants Preservation}: All physics constants were kept identical to ensure matching gameplay feel.
    \item \textbf{Rendering System Adaptation}: Java's Swing/AWT rendering was replaced with HTML5 Canvas.
    \item \textbf{Input System Rewrite}: Keyboard input handling was reimplemented for browser events.
\end{enumerate}

Key engine modules include:

\begin{lstlisting}
engine/
  MarioGame.ts       # Game loop (requestAnimationFrame)
  MarioWorld.ts      # World state and game logic
  MarioLevel.ts      # Level parsing from ASCII
  MarioSprite.ts     # Base sprite class
  sprites/           # Mario, enemies, items
  graphics/          # Camera, rendering
  effects/           # Visual effects (particles)
\end{lstlisting}

\subsection{Physics and Gameplay}
\label{subsec:physics}

The game engine preserves exact physics constants from the Java version to ensure identical gameplay:

\begin{table}[H]
\centering
\begin{tabular}{lr}
\toprule
\textbf{Parameter} & \textbf{Value} \\
\midrule
Ground Inertia & 0.89 \\
Air Inertia & 0.89 \\
Gravity & 3.0 \\
Jump Velocity & $-1.9$ \\
Walk Speed & 0.6 \\
Run Speed & 1.2 \\
\bottomrule
\end{tabular}
\caption{Core physics constants in the Mario engine.}
\label{tab:physics}
\end{table}

The game loop runs at 30 ticks per second (matching the original), with rendering at 60 FPS using \texttt{requestAnimationFrame}. Physics updates are deterministic, ensuring consistent gameplay across sessions.

\subsection{Level Format and Parsing}
\label{subsec:level-format}

Levels are represented as ASCII tilemaps, where each character corresponds to a tile type:

\begin{lstlisting}
--------------          (empty space)
----------S---          (S = breakable brick)
---E---Q------          (E = enemy, Q = question block)
##############          (# = ground)
\end{lstlisting}

The level parser (\texttt{MarioLevel.ts}) processes these tilemaps, converting characters to internal tile representations and spawning sprite entities for enemies and items. Validation rules ensure levels are playable:

\begin{itemize}
    \item Minimum dimensions: $50 \times 16$ tiles
    \item Maximum dimensions: $500 \times 16$ tiles
    \item Required: Mario spawn point and exit flag
    \item Prohibited: Unreachable exits, instant-death spawns
\end{itemize}

%==============================================================================
\section{Battle and Voting System}
\label{sec:voting}
%==============================================================================

\subsection{Battle Flow}
\label{subsec:battle-flow}

The core user experience of PCG Arena is the ``battle''---a blind A/B test where players compare two levels from different generators. The battle flow proceeds as follows:

\begin{enumerate}
    \item \textbf{Battle Request}: The frontend requests a new battle from \texttt{POST /v1/battles:next}.
    \item \textbf{Matchmaking}: The AGIS algorithm selects two generators and one random level from each.
    \item \textbf{Gameplay}: The player plays both levels sequentially (left then right).
    \item \textbf{Voting}: After playing both, the player votes for their preference.
    \item \textbf{Rating Update}: The backend updates Glicko-2 ratings atomically.
\end{enumerate}

Figure~\ref{fig:gameplay} shows the main gameplay interface where players experience levels side-by-side.

\begin{figure}[H]
    \centering
    \includegraphics[width=0.95\textwidth]{main_gameplay.png}
    \caption{Main gameplay view showing the active level with game controls and timer.}
    \label{fig:gameplay}
\end{figure}

\subsection{Voting Interface}
\label{subsec:voting-interface}

After completing both levels, players are presented with the voting interface (Figure~\ref{fig:voting}). The interface displays:

\begin{itemize}
    \item Static visualizations of both complete levels for reference
    \item Vote buttons: ``Left is Better,'' ``Right is Better,'' ``Tie,'' ``Skip''
    \item Optional tagging interface for additional feedback
\end{itemize}

\begin{figure}[H]
    \centering
    \includegraphics[width=0.9\textwidth]{voting_panel.png}
    \caption{Voting panel showing level previews and vote options after gameplay.}
    \label{fig:voting}
\end{figure}

The ``Skip'' option allows players to abstain from voting if they encountered technical issues or feel unable to judge. Skipped battles do not affect ratings.

\subsection{Tagging System}
\label{subsec:tagging}

In addition to the binary preference vote, players can optionally tag each level with descriptive labels. Tags capture qualitative aspects that the vote alone cannot express:

\begin{table}[H]
\centering
\begin{tabular}{ll}
\toprule
\textbf{Category} & \textbf{Tags} \\
\midrule
Difficulty & too-easy, just-right, too-hard \\
Design & creative, boring, frustrating, fun \\
Structure & linear, exploration, puzzle-like \\
Issues & broken, impossible, unplayable \\
\bottomrule
\end{tabular}
\caption{Available tags for level characterization.}
\label{tab:tags}
\end{table}

Tag data enables richer analysis, such as correlating perceived difficulty with win rates or identifying generators that produce consistently ``frustrating'' levels.

%==============================================================================
\section{Matchmaking Algorithm}
\label{sec:matchmaking}
%==============================================================================

\subsection{Problem Definition}
\label{subsec:matchmaking-problem}

The matchmaking problem in PCG Arena is: given $n$ generators with current ratings and uncertainties, select a pair for the next battle that maximizes information gain toward accurate final rankings.

Naive approaches fail:

\begin{itemize}
    \item \textbf{Uniform Random}: Wastes battles on already-converged pairs.
    \item \textbf{Most Uncertain First}: May neglect head-to-head comparisons between similar generators.
    \item \textbf{Round-Robin}: Ignores rating information; pairs dissimilar generators.
\end{itemize}

The ideal matchmaking algorithm should:

\begin{enumerate}
    \item Prioritize generators with high uncertainty (new or rarely played)
    \item Prefer matchups between generators of similar skill levels
    \item Ensure all pairs get sufficient battles for the confusion matrix
    \item Converge quickly to stable rankings
\end{enumerate}

\subsection{AGIS: Adaptive Glicko-Informed Selection}
\label{subsec:agis}

We developed \textbf{AGIS (Adaptive Glicko-Informed Selection)}, a matchmaking algorithm that uses Glicko-2 rating information to optimize battle selection. AGIS operates in two stages:

\textbf{Stage 1: First Generator Selection}

Each generator $i$ receives a selection weight $w_i$ based on:

\begin{equation}
w_i = (1 + \text{RD}_{\text{norm}, i})^2 \cdot \text{boost}(\text{games}_i) \cdot (1 + \alpha \cdot \text{quality}_i)
\end{equation}

where $\text{RD}_{\text{norm}}$ is the normalized rating deviation (high = uncertain), $\text{boost}$ amplifies weight for generators with few games, and $\text{quality}$ provides slight bias toward higher-rated generators after convergence.

\textbf{Stage 2: Second Generator Selection}

Given the first generator $g_1$, the second generator is selected with weight:

\begin{equation}
w_{1,j} = \alpha \cdot S(g_1, g_j) + \beta \cdot U(g_j) + \gamma \cdot C(g_1, g_j)
\end{equation}

where:
\begin{itemize}
    \item $S(g_1, g_j)$: Rating similarity (Gaussian decay from rating difference)
    \item $U(g_j)$: Uncertainty factor (based on $g_j$'s RD)
    \item $C(g_1, g_j)$: Coverage factor (inverse of pair's battle count)
\end{itemize}

\subsection{Generator Selection Weights}
\label{subsec:selection-weights}

The configurable parameters of AGIS (with defaults) are:

\begin{table}[H]
\centering
\begin{tabular}{llr}
\toprule
\textbf{Parameter} & \textbf{Description} & \textbf{Default} \\
\midrule
\texttt{AGIS\_MIN\_GAMES} & Games before ``converged'' & 30 \\
\texttt{AGIS\_TARGET\_BATTLES\_PER\_PAIR} & Minimum battles per pair & 10 \\
\texttt{AGIS\_RATING\_SIGMA} & Rating similarity std dev & 150.0 \\
\texttt{AGIS\_QUALITY\_BIAS} & Quality bias strength & 0.2 \\
\midrule
$\alpha$ & Rating similarity weight & 0.5 \\
$\beta$ & Uncertainty weight & 0.3 \\
$\gamma$ & Coverage weight & 0.2 \\
\bottomrule
\end{tabular}
\caption{AGIS algorithm parameters.}
\label{tab:agis-params}
\end{table}

%==============================================================================
\section{Rating System}
\label{sec:rating}
%==============================================================================

\subsection{Glicko-2 Algorithm}
\label{subsec:glicko2}

PCG Arena uses the Glicko-2 rating system~\cite{glickman2012}, which maintains three parameters for each generator:

\begin{itemize}
    \item \textbf{Rating ($\mu$)}: The skill estimate, displayed on a scale centered at 1000.
    \item \textbf{Rating Deviation (RD, $\phi$)}: The uncertainty in the rating. High RD means unreliable rating.
    \item \textbf{Volatility ($\sigma$)}: Expected fluctuation in performance over time.
\end{itemize}

The expected score when generator $A$ plays against generator $B$ is:

\begin{equation}
E(\mu_A, \mu_B, \phi_B) = \frac{1}{1 + \exp(-g(\phi_B)(\mu_A - \mu_B))}
\end{equation}

where $g(\phi) = \frac{1}{\sqrt{1 + 3\phi^2/\pi^2}}$ reduces the impact of uncertain opponents.

\subsection{Rating Updates}
\label{subsec:rating-updates}

After each battle, the winning generator receives a score of 1, the loser 0, and ties result in 0.5 for both. The rating update proceeds through several steps:

\begin{enumerate}
    \item Convert ratings to Glicko-2 scale (centered at 0, scaled by 173.7178)
    \item Compute variance $v$ based on opponent's RD
    \item Compute delta $\Delta$: actual minus expected score
    \item Update volatility $\sigma'$ using iterative algorithm
    \item Update RD: $\phi' = 1/\sqrt{1/\phi^2 + 1/v}$
    \item Update rating: $\mu' = \mu + \phi'^2 \cdot g(\phi_B) \cdot (s - E)$
    \item Convert back to display scale
\end{enumerate}

Rating updates are performed atomically within a database transaction to ensure consistency.

\subsection{Confidence Intervals}
\label{subsec:confidence}

A key advantage of Glicko-2 is that RD provides natural confidence intervals. The 95\% confidence interval for a generator's true rating is approximately:

\begin{equation}
\text{CI}_{95\%} = [\mu - 1.96 \cdot \text{RD}, \mu + 1.96 \cdot \text{RD}]
\end{equation}

This enables meaningful statements like ``Generator A is rated 1200 $\pm$ 50 with 95\% confidence,'' which is essential for research applications.

Rating bounds are enforced:
\begin{itemize}
    \item Minimum rating: 100 (prevents extreme negative outliers)
    \item Maximum rating: 3000 (theoretical maximum)
    \item Minimum RD: 30 (very confident)
    \item Maximum RD: 350 (new/inactive)
\end{itemize}

%==============================================================================
\section{Leaderboard and Analytics}
\label{sec:leaderboard}
%==============================================================================

\subsection{Leaderboard Display}
\label{subsec:leaderboard-display}

The public leaderboard (Figure~\ref{fig:leaderboard}) displays all active generators ranked by their Glicko-2 rating. For each generator, the display shows:

\begin{itemize}
    \item Rank position
    \item Generator name and version
    \item Current rating with $\pm$ RD uncertainty
    \item Number of battles played
    \item Win/loss/tie record
\end{itemize}

\begin{figure}[H]
    \centering
    \includegraphics[width=0.95\textwidth]{leaderboard.png}
    \caption{Public leaderboard showing generator rankings with Glicko-2 ratings and uncertainty intervals.}
    \label{fig:leaderboard}
\end{figure}

The leaderboard updates in real-time as votes are submitted, allowing users to see the immediate impact of their participation.

\subsection{Generator Statistics}
\label{subsec:generator-stats}

Clicking on a generator in the leaderboard reveals detailed statistics (Figure~\ref{fig:generator-preview}):

\begin{itemize}
    \item Rating history over time
    \item Win rate against each opponent
    \item Level gallery with performance metrics per level
    \item Tag distribution (what players say about this generator's levels)
    \item Average play statistics (deaths, completion rate, time)
\end{itemize}

\begin{figure}[H]
    \centering
    \includegraphics[width=0.95\textwidth]{generator_preview.png}
    \caption{Generator detail page showing level gallery and statistics.}
    \label{fig:generator-preview}
\end{figure}

\subsection{Confusion Matrix}
\label{subsec:confusion-matrix}

The confusion matrix (Figure~\ref{fig:confusion}) provides a comprehensive view of head-to-head performance between all generator pairs. Each cell shows the win rate when the row generator faces the column generator.

\begin{figure}[H]
    \centering
    \includegraphics[width=0.8\textwidth]{confusion_matrix.png}
    \caption{Confusion matrix showing pairwise win rates between generators.}
    \label{fig:confusion}
\end{figure}

The matrix is particularly useful for:
\begin{itemize}
    \item Identifying rock-paper-scissors relationships (A beats B, B beats C, C beats A)
    \item Detecting generators with inconsistent performance
    \item Ensuring sufficient coverage---empty or low-count cells indicate gaps
\end{itemize}

Coverage gaps (Figure~\ref{fig:coverage}) are highlighted to guide future data collection.

\begin{figure}[H]
    \centering
    \includegraphics[width=0.8\textwidth]{coverage_gaps.png}
    \caption{Coverage analysis showing pairs with insufficient battle data.}
    \label{fig:coverage}
\end{figure}

%==============================================================================
\section{Authentication and User Management}
\label{sec:auth}
%==============================================================================

\subsection{Google OAuth Integration}
\label{subsec:google-oauth}

PCG Arena supports Google OAuth 2.0 for quick sign-in, leveraging Google's secure identity verification. The flow proceeds as:

\begin{enumerate}
    \item User clicks ``Sign In with Google''
    \item Browser opens Google OAuth popup
    \item User authenticates with Google credentials
    \item Google returns an ID token (JWT) to the frontend
    \item Frontend sends token to \texttt{POST /v1/auth/google}
    \item Backend verifies token with Google's servers
    \item Backend creates or retrieves user record
    \item Backend issues session cookie (\texttt{arena\_session})
    \item User is redirected to Builder Profile
\end{enumerate}

Google OAuth users are automatically marked as email-verified (Google handles this), enabling immediate access to submission features.

\subsection{Email/Password Authentication}
\label{subsec:email-auth}

For users who prefer not to use Google, PCG Arena supports traditional email/password authentication:

\textbf{Registration:}
\begin{enumerate}
    \item User provides email, password, and display name
    \item Backend validates email format and password strength
    \item Password is hashed using bcrypt (work factor 12)
    \item User record created with \texttt{is\_email\_verified = 0}
    \item Verification email sent (see Section~\ref{subsec:email-verification})
\end{enumerate}

\textbf{Login:}
\begin{enumerate}
    \item User provides email and password
    \item Backend retrieves user record by email
    \item Password verified against stored hash
    \item Session cookie issued on successful verification
\end{enumerate}

Security measures include:
\begin{itemize}
    \item Password minimum length: 8 characters
    \item Session tokens: 256-bit cryptographically random
    \item Session expiry: 7 days (configurable)
    \item HTTP-only cookies prevent XSS token theft
\end{itemize}

\subsection{Email Verification with Twilio SendGrid}
\label{subsec:email-verification}

Email verification ensures that users control the email addresses they register with. PCG Arena uses Twilio SendGrid for email delivery:

\begin{enumerate}
    \item On registration, generate a random verification token (64 hex characters)
    \item Store token in \texttt{email\_verifications} table with 24-hour expiry
    \item Send verification email via SendGrid API containing link: \\ \texttt{https://pcg-arena.com/verify-email?token=<token>}
    \item User clicks link, frontend calls \texttt{POST /v1/auth/verify-email}
    \item Backend verifies token validity and expiry
    \item User marked as verified; token deleted
\end{enumerate}

Password reset follows a similar flow, using separate tokens stored in \texttt{password\_resets}. Reset tokens expire after 1 hour for security.

%==============================================================================
\section{Builder Profile}
\label{sec:builder}
%==============================================================================

\subsection{Generator Submission}
\label{subsec:generator-submission}

The Builder Profile (Figure~\ref{fig:builder}) enables authenticated researchers to submit their own generators to compete on the leaderboard. The submission process:

\begin{enumerate}
    \item User navigates to Builder Profile (requires authentication)
    \item Fills out generator metadata: name, version, description, tags, URL
    \item Uploads a ZIP file containing level files (ASCII tilemaps)
    \item Backend validates all levels (see Section~\ref{subsec:level-validation})
    \item Generator and levels are inserted into the database
    \item Initial Glicko-2 rating assigned (1000 $\pm$ 350)
    \item Generator appears on leaderboard immediately
\end{enumerate}

\begin{figure}[H]
    \centering
    \includegraphics[width=0.95\textwidth]{builder_profile.png}
    \caption{Builder Profile page for managing submitted generators.}
    \label{fig:builder}
\end{figure}

Constraints on submissions:
\begin{itemize}
    \item Maximum 3 generators per user (prevents flooding)
    \item 50--200 levels per generator (ensures variety without overwhelming)
    \item Each level must pass validation checks
\end{itemize}

\subsection{Level Validation}
\label{subsec:level-validation}

Every submitted level undergoes validation to ensure playability:

\begin{enumerate}
    \item \textbf{Format Check}: File must be valid ASCII, correct dimensions
    \item \textbf{Character Validation}: Only recognized tile characters allowed
    \item \textbf{Spawn/Exit Check}: Level must have Mario spawn point and reachable exit
    \item \textbf{Deduplication}: Content hash computed; duplicates rejected
    \item \textbf{Playability Heuristics}: Basic reachability analysis
\end{enumerate}

Validation errors return detailed feedback indicating which levels failed and why, enabling researchers to fix issues before resubmitting.

\subsection{Generator Management}
\label{subsec:generator-management}

After submission, researchers can manage their generators:

\textbf{Update Version:}
\begin{itemize}
    \item Upload new level set while preserving rating
    \item Version number incremented
    \item Old levels replaced with new ones
    \item Rating history maintained for continuity
\end{itemize}

\textbf{Delete Generator:}
\begin{itemize}
    \item Generators with battles: soft delete (marked inactive, preserved for historical data)
    \item Generators without battles: hard delete (completely removed)
    \item Soft-deleted generators don't appear in matchmaking but remain in historical rankings
\end{itemize}

This design ensures data integrity---votes referencing deleted generators remain valid for research analysis.

%==============================================================================
\section{Admin Dashboard}
\label{sec:admin}
%==============================================================================

\subsection{Platform Overview}
\label{subsec:admin-overview}

The Admin Dashboard (Figure~\ref{fig:admin-dashboard}) provides platform administrators with comprehensive oversight and management capabilities. Access is restricted to users whose email addresses are listed in the \texttt{ARENA\_ADMIN\_EMAILS} configuration.

\begin{figure}[H]
    \centering
    \includegraphics[width=0.95\textwidth]{admin_dashboard.png}
    \caption{Admin Dashboard showing platform statistics and management tools.}
    \label{fig:admin-dashboard}
\end{figure}

Key metrics displayed:
\begin{itemize}
    \item Total battles completed and pending
    \item Total votes and vote distribution (left/right/tie/skip)
    \item Active generators and users
    \item Recent activity timeline
    \item System health indicators
\end{itemize}

Administrative actions available:
\begin{itemize}
    \item View and ban builder accounts
    \item Delete inappropriate generators
    \item Monitor matchmaking statistics
    \item Trigger manual data exports
\end{itemize}

\subsection{Data Export}
\label{subsec:data-export}

PCG Arena provides comprehensive data export for research purposes (Figure~\ref{fig:data-export}). Export endpoints (admin-only) include:

\begin{table}[H]
\centering
\begin{tabular}{ll}
\toprule
\textbf{Endpoint} & \textbf{Data} \\
\midrule
\texttt{/v1/admin/export/votes} & All votes with telemetry and tags \\
\texttt{/v1/admin/export/trajectories} & Player position traces \\
\texttt{/v1/admin/export/level-stats} & Aggregate per-level metrics \\
\texttt{/v1/admin/export/player-profiles} & Anonymous player IDs and sessions \\
\bottomrule
\end{tabular}
\caption{Available data export endpoints.}
\label{tab:exports}
\end{table}

\begin{figure}[H]
    \centering
    \includegraphics[width=0.9\textwidth]{data_export.png}
    \caption{Data export interface for downloading research datasets.}
    \label{fig:data-export}
\end{figure}

Data is exported in JSON and CSV formats suitable for analysis with standard tools (Python/pandas, R, etc.).

\subsection{Research Analytics}
\label{subsec:research-analytics}

The platform collects rich telemetry enabling advanced research analytics:

\textbf{Level Statistics} (Figure~\ref{fig:statistics}):
\begin{itemize}
    \item Completion rate: percentage of players who finish
    \item One-attempt failure/death rate
    \item Average play time
    \item Win rate when featured in battles
\end{itemize}

\begin{figure}[H]
    \centering
    \includegraphics[width=0.95\textwidth]{statistics_page.png}
    \caption{Statistics page showing detailed level and generator analytics.}
    \label{fig:statistics}
\end{figure}

\textbf{Death Heatmaps}:
\begin{itemize}
    \item Aggregate death locations across all plays of a level
    \item Reveals difficulty spikes and design issues
    \item Exportable for external visualization
\end{itemize}

\textbf{Detailed Level View} (Figure~\ref{fig:detailed-level}):
\begin{itemize}
    \item Full level visualization with overlay options
    \item Play trajectory visualization
    \item Per-tile analysis
\end{itemize}

\begin{figure}[H]
    \centering
    \includegraphics[width=0.95\textwidth]{detailed_level_view.png}
    \caption{Detailed level view with analytics overlays.}
    \label{fig:detailed-level}
\end{figure}

%==============================================================================
\section{Deployment and Infrastructure}
\label{sec:deployment}
%==============================================================================

\subsection{Google Cloud Platform Setup}
\label{subsec:gcp}

PCG Arena is deployed on Google Cloud Platform (GCP) using a single Compute Engine virtual machine. This setup balances cost-effectiveness with reliability for a research platform:

\textbf{VM Specifications:}
\begin{itemize}
    \item Machine type: \texttt{e2-micro} (free tier eligible)
    \item Region: \texttt{us-central1} (low latency for most users)
    \item Boot disk: Ubuntu 22.04 LTS, 10 GB standard persistent disk
    \item Static external IP: reserved for DNS stability
\end{itemize}

\textbf{Firewall Configuration:}
\begin{itemize}
    \item Port 80 (HTTP): Allowed for redirect to HTTPS
    \item Port 443 (HTTPS): Allowed for all traffic
    \item Port 8080 (backend): \textbf{NOT} public; backend binds to localhost
\end{itemize}

The architecture ensures that all public traffic goes through the Caddy reverse proxy, which handles TLS termination and routes requests appropriately.

\subsection{Docker Containerization}
\label{subsec:docker}

The backend runs in a Docker container, providing:

\begin{itemize}
    \item Consistent environment across development and production
    \item Easy deployment via \texttt{docker compose up -d --build}
    \item Isolated dependencies (Python packages)
    \item Simple rollback by switching image tags
\end{itemize}

The \texttt{docker-compose.yml} configuration:

\begin{lstlisting}[language=bash]
services:
  backend:
    build: ./backend
    ports:
      - "127.0.0.1:8080:8080"  # Localhost only
    volumes:
      - ./db/local:/data
      - ./db/migrations:/migrations
      - ./db/seed:/seed
    environment:
      - ARENA_DB_PATH=/data/arena.sqlite
\end{lstlisting}

Database files are bind-mounted from the host, ensuring persistence across container rebuilds.

\subsection{Domain and HTTPS Configuration}
\label{subsec:domain}

The platform is accessible at \url{https://pcg-arena.com}, configured as follows:

\textbf{Domain Registration:}
\begin{itemize}
    \item Domain: \texttt{pcg-arena.com} (registered via Namecheap)
    \item DNS A records: \texttt{@} and \texttt{www} point to VM's static IP
\end{itemize}

\textbf{HTTPS via Caddy:}
\begin{itemize}
    \item Automatic TLS certificate provisioning (Let's Encrypt)
    \item Automatic certificate renewal
    \item HTTP to HTTPS redirect
    \item Reverse proxy to backend (\texttt{localhost:8080})
    \item Static file serving for frontend build
\end{itemize}

Caddyfile configuration:

\begin{lstlisting}
pcg-arena.com, www.pcg-arena.com {
    root * /var/www/pcg-arena/frontend/dist
    file_server
    
    @api path /v1/* /health
    handle @api {
        reverse_proxy localhost:8080
    }
    
    try_files {path} /index.html
}
\end{lstlisting}

\subsection{Monitoring and Alerts}
\label{subsec:monitoring}

GCP provides built-in monitoring capabilities:

\begin{itemize}
    \item \textbf{Uptime Checks}: HTTP(S) probe to \texttt{/health} endpoint every 60 seconds
    \item \textbf{Alerting Policies}: Email notification when health check fails
    \item \textbf{VM Metrics}: CPU, memory, disk usage monitoring
    \item \textbf{Logging}: Application logs accessible via Cloud Console
\end{itemize}

Automated backups run daily via cron job:
\begin{lstlisting}[language=bash]
0 3 * * * /home/user/pcg-arena/backend/scripts/backup.sh
\end{lstlisting}

Backups are stored in \texttt{/backups/} with 7-day retention, enabling recovery from data corruption or accidental deletion.

%==============================================================================
\section{Results and Discussion}
\label{sec:results}
%==============================================================================

This section reframes the platform data as an exploratory user study and generator characterisation. Rather than testing a long list of isolated hypotheses, the analysis asks five broader research questions:

\begin{enumerate}
    \item \textbf{RQ1}: Which generators are preferred by players in blind pairwise comparisons?
    \item \textbf{RQ2}: How do generators differ in static expressive range, and which structural properties align with player preference?
    \item \textbf{RQ3}: What behavioral fingerprints appear in player trajectories for each generator?
    \item \textbf{RQ4}: Are anonymous player preferences and play styles homogeneous, or do visible subgroups emerge?
    \item \textbf{RQ5}: What do optional tags reveal about why levels win or lose?
\end{enumerate}

The emphasis is descriptive and visual: the goal is to show what PCG Arena can measure---preference, structure, gameplay behavior, user heterogeneity, and qualitative failure modes---without over-claiming universal design laws from a sparse exploratory dataset.

\subsection{Study Protocol and Dataset}
\label{subsec:study-protocol-dataset}

Participants were recruited through university social-media channels, word of mouth among game-development and computer-science communities, and a small promoted Reddit campaign targeting game-development and procedural-generation communities. Participation was anonymous. The platform assigned browser-based anonymous player identifiers and session identifiers, persisted through cookies or local browser storage. Consequently, ``players'' in this analysis should be interpreted as anonymous browser/player IDs rather than verified unique humans; clearing cookies or changing device/browser could create a new ID.

Each battle used a blind A/B protocol. A player saw two levels from different generators, played one attempt on each side, and then selected the preferred side, a tie, or skip. Optional tags could be attached to either side. The one-attempt design is important for interpretation: the death variable is a binary failure indicator for a side-level attempt, not a repeated-death count from multiple retries.

The May 2026 analysis export contains the data summarized in Table~\ref{tab:eda-dataset-summary}. The research plots exclude the diagnostic \texttt{test-gen} generator and therefore compare 14 substantive generator IDs. Skip votes are included in engagement summaries but excluded from preference rankings. Ties are treated as 0.5 points for each side.

\begin{table}[H]
\centering
\caption{Dataset summary for the exploratory analysis. Player counts refer to anonymous browser/player IDs, not verified unique humans.}
\label{tab:eda-dataset-summary}
\begin{tabular}{lr}
\toprule
\textbf{Quantity} & \textbf{Count} \\
\midrule
Vote rows analyzed & 1,000 \\
Anonymous player profiles & 79 \\
Anonymous player IDs in vote table & 77 \\
Session IDs in vote table & 213 \\
Level-stat records exported & 1,104 \\
Generator IDs in research plots & 14 \\
Per-side telemetry records & 2,000 \\
Played side records & 1,968 \\
Embedded trajectory records & 1,968 \\
Tag assignments & 329 \\
\bottomrule
\end{tabular}
\end{table}

The vote outcomes are balanced between the two screen sides (418 left wins and 398 right wins), with 165 ties and 19 skips. The median play duration is 8.0 seconds, but a small number of inactive-tab outliers are present; duration-dependent plots therefore use clipping or robust summaries. Figure~\ref{fig:eda-study-overview} summarizes the analysis links among anonymous players, votes, side-level telemetry, tags, trajectories, and static level text.

\begin{figure}[H]
\centering
\includegraphics[width=0.95\textwidth]{eda_study_overview.png}
\caption{Analysis data model for the revised EDA. Pairwise votes connect anonymous player/session IDs to side-level telemetry, optional tags, embedded trajectories, and generator/level identifiers.}
\label{fig:eda-study-overview}
\end{figure}

\subsection{RQ1: Generator Preference Ranking}
\label{subsec:eda-generator-preference}

The core purpose of PCG Arena is to measure human preference between generators. For each non-skip battle, the selected side receives score 1, the rejected side receives score 0, and ties give 0.5 to both sides. Figure~\ref{fig:eda-generator-ranking} shows the resulting generator leaderboard with bootstrap uncertainty intervals over votes.

\begin{figure}[H]
\centering
\includegraphics[width=0.95\textwidth]{eda_generator_ranking.png}
\caption{Generator ranking from blind pairwise votes. Bars show mean preference score, where win = 1, tie = 0.5, and loss = 0; horizontal intervals show bootstrap uncertainty. The text at the end of each bar gives the win-loss-tie record.}
\label{fig:eda-generator-ranking}
\end{figure}

The strongest result is that the human-authored \texttt{original} baseline and \texttt{mariodpo} are nearly tied at the top. \texttt{original} scores 0.833, while \texttt{mariodpo} scores 0.832. Both are well above the neutral 0.5 line. The next tier consists of \texttt{ore} (0.655) and \texttt{mariogpt} (0.643), followed by mid-table constructive and neural generators. The pattern-based generators occupy the bottom of the ranking: \texttt{patternWeightCount} scores 0.385, \texttt{patternOccur} scores 0.284, and \texttt{patternCount} scores 0.262.

This ordering is meaningful for two reasons. First, the human-authored baseline remains a strong benchmark, which is expected if the vote data captures level quality. Second, \texttt{mariodpo} approaches the human baseline and clearly outperforms the other procedural generators in this export, suggesting that preference-aligned generation is a promising direction for the platform.

Figure~\ref{fig:eda-pairwise-confusion} gives the corresponding head-to-head matrix. Rows and columns are ordered by the ranking above. The matrix reveals that the top generators are not merely benefiting from weak opponents: \texttt{original} and \texttt{mariodpo} perform well across many pairings, whereas pattern generators lose consistently across opponents.

\begin{figure}[H]
\centering
\includegraphics[width=0.95\textwidth]{eda_pairwise_confusion.png}
\caption{Pairwise generator preference matrix ordered by overall ranking. Each cell is the score rate for the row generator against the column generator, with ties counted as 0.5. Lighter green indicates a higher score rate for the row generator.}
\label{fig:eda-pairwise-confusion}
\end{figure}

\subsection{RQ2: Static Expressive Range}
\label{subsec:eda-static-expressive-range}

The exported aggregate \texttt{level\_stats} table does not yet contain populated structural columns, so structural metrics were computed directly from the ASCII level text files. This separates generator characterization from gameplay outcomes. The static corpus is larger than the live vote export for some generators, especially \texttt{mariodpo}, and should be interpreted as a generator-level design-space sample rather than a count of voted levels.

The metrics are operational definitions chosen for consistency across generators. Linearity is the $R^2$ of a least-squares line fit to the detected playable surface profile. Density metrics count solid, enemy, reward, and gap structure per tile or per level column. Leniency combines gap density, enemy density, maximum gap width, rewards, and blocks into a single descriptive score; because leniency definitions vary in the PCG literature, it is used here as an operational descriptor rather than a universal measure of difficulty. Diversity is estimated with tile/column entropy, unique-column ratio, compression ratio, and gzip normalized compression distance (NCD).

Figure~\ref{fig:eda-static-metrics-table} summarizes the main static metrics by generator. Figure~\ref{fig:eda-expressive-range} then shows each generator's expressive range in the plane of linearity and operational leniency.

\begin{figure}[H]
\centering
\includegraphics[width=0.95\textwidth]{eda_static_metrics_table.png}
\caption{Static expressive metrics computed from level text files. Metrics include fitted-surface linearity, operational leniency, gap density, enemy density, and within-generator compression distance.}
\label{fig:eda-static-metrics-table}
\end{figure}

\begin{figure}[H]
\centering
\includegraphics[width=\textwidth]{eda_expressive_range.png}
\caption{Expressive-range view of static level structure. Each subplot shows levels from one generator in the linearity--leniency plane, ordered by preference ranking.}
\label{fig:eda-expressive-range}
\end{figure}

The static analysis shows that generators occupy distinct regions of design space. Pattern generators have zero reward density in the parsed corpus and comparatively high enemy density, which helps explain why they are interpretable as harsh or unrewarding despite sometimes receiving high operational leniency scores from non-gap structure alone. \texttt{mariogpt} and \texttt{original} have the highest reward densities among the compared generators. \texttt{ore} has the highest unique-column ratio and high within-generator NCD, suggesting broader structural diversity.

Figure~\ref{fig:eda-static-vs-rating} relates selected structural summaries to preference score. With only 14 generators, these plots are best read as effect-size and direction summaries rather than formal hypothesis tests. Reward density and unique-column ratio show strong positive rank associations with preference score in this export, while distance from the \texttt{original} corpus by NCD is negatively associated with preference. Enemy density alone is not sufficient to explain rankings, reinforcing that generator quality is multi-dimensional.

\begin{figure}[H]
\centering
\includegraphics[width=\textwidth]{eda_static_metrics_vs_rating.png}
\caption{Generator preference score against selected static metrics. The fitted lines are descriptive trend guides, not causal claims.}
\label{fig:eda-static-vs-rating}
\end{figure}

\subsection{RQ3: Gameplay Trajectory Fingerprints}
\label{subsec:eda-trajectory-fingerprints}

Trajectory recording was available in embedded vote telemetry for nearly all played side records in this export. For each generator, Figure~\ref{fig:eda-trajectory-stacks} overlays all recorded Mario paths as thin red traces. This plot is intentionally simple: it shows where players actually moved in levels from each generator under the one-attempt battle protocol.

\begin{figure}[H]
\centering
\includegraphics[width=\textwidth]{eda_trajectory_stacks.png}
\caption{Stacked gameplay trajectories by generator. Each subplot overlays all recorded trajectories for one generator as thin red traces on a white background, ordered by preference score. The vertical axis is inverted to match screen coordinates.}
\label{fig:eda-trajectory-stacks}
\end{figure}

The trajectory stacks make generator differences visually apparent. Stronger generators such as \texttt{original} and \texttt{mariodpo} show paths that often progress substantially through the level and occupy a broad portion of the play space. Several low-ranked pattern generators show many traces terminating early or concentrating in narrow regions, consistent with the high early-failure rates computed from death locations.

Figure~\ref{fig:eda-trajectory-metrics} reduces these visual patterns to generator-level summaries. Occupancy entropy measures how broadly trajectories cover a fixed spatial grid. Median maximum $x$ measures progress through the level. Verticality is the mean within-trajectory standard deviation of $y$. Death concentration is the largest share of death locations falling into a single horizontal bin. These are behavioral descriptors of a one-attempt protocol: death-location metrics indicate where failures happen, not how often players retry the same level.

\begin{figure}[H]
\centering
\includegraphics[width=0.95\textwidth]{eda_trajectory_metrics_vs_rating.png}
\caption{Trajectory-derived behavioral summaries versus generator preference score. Higher occupancy entropy and lower death concentration are associated with stronger generators in this export.}
\label{fig:eda-trajectory-metrics}
\end{figure}

The strongest trajectory-level patterns are occupancy entropy and death concentration. Generator-level occupancy entropy has a positive rank association with preference score, while concentrated death locations have a negative association. The weakest pattern generator, \texttt{patternCount}, has median maximum progress of only 174.5 pixels and a max death-bin share of 0.63, meaning failures are highly concentrated early in the level. In contrast, \texttt{original} and \texttt{mariodpo} have median progress around 1126.5 and 988.0 pixels respectively, with lower death concentration.

The deployed analytics UI also supports level-level inspection beyond aggregate plots. Figure~\ref{fig:eda-detailed-level-stats} shows the Detailed Level Statistics page, which combines the level view, player traces, death locations, and tag counts. This is important as a platform contribution: PCG Arena is not only a leaderboard, but also an instrument for diagnosing why a generator wins or loses.

\begin{figure}[H]
\centering
\includegraphics[width=0.95\textwidth]{detailed_level_statistics_page.png}
\caption{Detailed Level Statistics page from PCG Arena. The interface supports qualitative inspection of individual levels through trajectory overlays, death locations, tag counts, and level visualization.}
\label{fig:eda-detailed-level-stats}
\end{figure}

\subsection{RQ4: User Preference Heterogeneity}
\label{subsec:eda-user-heterogeneity}

The anonymous-player data are sparse and should not be treated as demographic information, but they are still useful for understanding engagement and possible preference heterogeneity. Figure~\ref{fig:eda-engagement} shows votes over time, votes per anonymous player ID, and cumulative vote contribution.

\begin{figure}[H]
\centering
\includegraphics[width=\textwidth]{eda_engagement.png}
\caption{Engagement concentration in the anonymous player-ID data. The analysis uses browser/player IDs, so these counts should not be interpreted as verified unique humans.}
\label{fig:eda-engagement}
\end{figure}

The vote distribution is highly concentrated. Among the 77 anonymous player IDs appearing in the vote table, the median player ID contributes 6 votes, while the most active ID contributes 217 votes. The top 5 IDs account for 54\% of the analyzed votes, and the top 10 account for 63.5\%. This concentration is typical for voluntary web studies, but it matters for interpretation: aggregate generator scores reflect both broad participation and the preferences of a small number of highly engaged player IDs.

Figure~\ref{fig:eda-user-heatmap} examines the player-by-generator preference matrix for anonymous player IDs with at least five votes. Each cell is the mean score assigned by a player ID when that generator appeared. Missing player-generator combinations are filled with the neutral value 0.5 only for visualization and clustering.

\begin{figure}[H]
\centering
\includegraphics[width=\textwidth]{eda_user_generator_heatmap.png}
\caption{Anonymous player ID by generator preference heatmap for player IDs with at least five votes. Rows are hierarchically clustered and columns follow the global generator ranking. The plot is exploratory because many player-generator combinations are sparse.}
\label{fig:eda-user-heatmap}
\end{figure}

The heatmap suggests that the player population should not be assumed homogeneous. Many player IDs follow the global ordering by favoring \texttt{original}, \texttt{mariodpo}, or other high-ranked generators, but some rows show more tolerant or idiosyncratic preferences for mid-table generators. Because IDs are anonymous and sparse, this is best interpreted as exploratory evidence for preference diversity rather than a definitive taxonomy of player types.

\subsection{RQ5: Tag Semantics and Failure Modes}
\label{subsec:eda-tag-semantics}

Tags provide a lightweight qualitative layer over the binary vote. They are optional, sparse, and subjective, so they should not be treated as ground-truth labels. Their value is diagnostic: they help distinguish why a level wins or loses.

Figure~\ref{fig:eda-tag-semantics} summarizes tag frequency, tag rate by generator, tag co-occurrence, and outcome distribution when each tag is present. The most common tags are \texttt{impossible} (87 assignments), \texttt{fun} (56), \texttt{too\_hard} (55), \texttt{boring} (49), \texttt{creative} (39), \texttt{broken\_graphics} (22), and \texttt{too\_easy} (21). Several available tags were unused in this export.

\begin{figure}[H]
\centering
\includegraphics[width=\textwidth]{eda_tag_semantics.png}
\caption{Tag semantics summary. The panels show tag frequency, tag rate by generator, tag co-occurrence, and the outcome distribution for tagged sides.}
\label{fig:eda-tag-semantics}
\end{figure}

Positive tags are strongly aligned with preference outcomes. Sides tagged \texttt{creative} win 87\% of the time and sides tagged \texttt{fun} win 71\% of the time. Negative tags behave differently: \texttt{impossible}, \texttt{too\_hard}, and \texttt{boring} are much more common on losing sides than winning sides. The tag-by-generator heatmap also makes generator failure modes interpretable: some generators lose because they are marked as impossible or too hard, while others are rejected as boring or visually broken.

\subsection{Summary of Findings}
\label{subsec:eda-findings-summary}

The revised analysis supports a stronger story than a sequence of hypothesis tests. PCG Arena provides a live human-preference measurement instrument for procedural level generators. It ranks generators through blind pairwise comparison, shows stable head-to-head patterns, connects rankings to static expressive range, reveals behavioral differences through trajectories, exposes engagement concentration and possible preference diversity, and uses tags to explain qualitative failure modes.

The main empirical findings are:

\begin{enumerate}
    \item \textbf{Preference ranking}: \texttt{original} and \texttt{mariodpo} are the top generators in the analyzed export, with nearly identical preference scores around 0.83.
    \item \textbf{Generator families}: Pattern-based generators are consistently weak, while neural and constructive generators occupy the middle tiers.
    \item \textbf{Static structure}: Generators occupy visibly distinct expressive ranges. Reward density, unique-column ratio, and distance to the original corpus are more informative than any single difficulty proxy.
    \item \textbf{Gameplay behavior}: Trajectories reveal navigability and failure modes. Stronger generators tend to support broader movement footprints and less concentrated failure locations.
    \item \textbf{Users}: Vote contribution is concentrated among a small number of highly active anonymous player IDs, and the user-generator heatmap suggests exploratory preference heterogeneity.
    \item \textbf{Tags}: Optional tags add semantic explanations to votes, distinguishing positive qualities such as \texttt{fun} and \texttt{creative} from failure modes such as \texttt{impossible}, \texttt{too\_hard}, and \texttt{boring}.
\end{enumerate}

These findings motivate the next section: if human preference data can reliably identify high-quality generators and interpretable failure modes, it can also be used to train and evaluate preference-aligned generators such as Mario-DPO.




%==============================================================================
\section{Mario-DPO: Toward SOTA Level Generation}
\label{sec:mario-dpo}
%==============================================================================

Building on the preference ranking, expressive-range analysis, trajectory fingerprints, and tag diagnostics above, this section presents \textbf{Mario-DPO}, a proposed framework for training a state-of-the-art procedural level generator using Direct Preference Optimization (DPO)~\cite{rafailov2023dpo}. The goal is to create a generator that achieves statistical parity ($>50\%$ win rate) against human-designed ``Original'' levels in PCG Arena.

\subsection{Framework Architecture}
\label{subsec:dpo-architecture}

Mario-DPO follows a three-phase training pipeline:

\textbf{Phase 1: Supervised Fine-Tuning (SFT)}
\begin{itemize}
    \item Base model: GPT-2 (124M parameters) with character-level tokenization
    \item Training corpus: 15 Original (Nintendo-style) levels from the seed dataset
    \item Token vocabulary: 37 unique ASCII tile characters
    \item Context length: 2048 tokens ($\approx$ 128 columns $\times$ 16 rows)
\end{itemize}

\textbf{Phase 2: Preference-Data Expansion}

The revised EDA suggests that an automatic judge should be treated as a diagnostic prior rather than a replacement for human votes. A practical judge can combine static expressive-range features with trajectory-derived behavior summaries:
\begin{equation}
J_{final} = J_{static} + J_{behavior}
\end{equation}
where:
\begin{equation}
J_{static} = w_{style} \cdot \frac{1}{1 + D_{orig}} + w_{reward} \cdot \text{RewardDensity} - w_{gap} \cdot \text{GapDensity} - w_{enemy} \cdot \text{EnemyDensity}
\end{equation}
\begin{equation}
J_{behavior} = w_{progress} \cdot \text{Progress} + w_{entropy} \cdot \text{OccupancyEntropy} - w_{death} \cdot \text{DeathConcentration}
\end{equation}

Here $D_{orig}$ denotes distance from the human-authored \texttt{original} corpus in the static feature space. The weights should be fitted or calibrated against PCG Arena vote data, with human votes weighted more strongly than synthetic preferences.

\textbf{Phase 3: DPO Fine-Tuning}

Apply Direct Preference Optimization using the combined preference dataset:
\begin{equation}
\mathcal{L}_{DPO}(\pi_\theta; \pi_{ref}) = -\mathbb{E}_{(x, y_w, y_l) \sim D} \left[ \log \sigma \left( \beta \log \frac{\pi_\theta(y_w|x)}{\pi_{ref}(y_w|x)} - \beta \log \frac{\pi_\theta(y_l|x)}{\pi_{ref}(y_l|x)} \right) \right]
\end{equation}

\subsection{Preference Dataset Analysis}
\label{subsec:preference-dataset}

The side-level representation from Section~\ref{sec:results} can be transformed directly into DPO preference pairs. Decisive votes provide winner/loser pairs, ties can be retained as weak or neutral preferences, and skips are excluded from training. In the analyzed export, 816 votes are decisive, 165 are ties, and 19 are skips.

\begin{table}[H]
\centering
\caption{Preference data available for DPO-style training from the analyzed export. Human pairs should receive higher weight than any synthetic judge-generated pairs.}
\label{tab:dpo-dataset}
\begin{tabular}{llll}
\toprule
\textbf{Source} & \textbf{Pairs} & \textbf{Suggested Use} & \textbf{Weight} \\
\midrule
Human decisive votes & 816 & winner/loser DPO pairs & high \\
Human ties & 165 & weak or neutral comparisons & optional \\
Skipped votes & 19 & exclude from preference training & none \\
Synthetic judge pairs & configurable & coverage expansion & low \\
\bottomrule
\end{tabular}
\end{table}

The generator ranking in Figure~\ref{fig:eda-generator-ranking} provides the main validation target for any synthetic judge. A useful judge should preserve the broad human ordering: \texttt{original} and \texttt{mariodpo} near the top, competitive neural and constructive generators in the middle, and pattern-based generators near the bottom. Before using synthetic pairs at scale, the judge should be calibrated against held-out human votes and inspected for failure modes using the tag and trajectory diagnostics above.

Key implications for Mario-DPO are:
\begin{itemize}
    \item \textbf{Human preferences remain primary}: real votes define the target ranking and should dominate the loss.
    \item \textbf{Static features are useful filters}: gap, reward, enemy, and style-distance metrics can cheaply reject poor samples before human evaluation.
    \item \textbf{Trajectory features expose navigability}: progress, occupancy entropy, and death concentration help detect levels that look valid but play poorly.
\end{itemize}

\subsection{Level Corpus Statistics}
\label{subsec:corpus-stats}

Analysis of the complete level corpus:

\begin{table}[H]
\centering
\caption{Level Corpus Summary}
\label{tab:corpus-summary}
\begin{tabular}{lr}
\toprule
\textbf{Metric} & \textbf{Value} \\
\midrule
Total Level Files & 1,215 \\
Original (Training) Levels & 15 \\
Tile Vocabulary Size & 37 \\
Average Level Height & 15.7 tiles \\
Average Level Width & 193.8 tiles \\
Levels per Generator & $\sim$100 \\
\bottomrule
\end{tabular}
\end{table}

The 37-character ASCII tile vocabulary enables efficient character-level tokenization without subword complexity.

\subsection{Proposed Training Configuration}
\label{subsec:training-config}

Based on the dataset analysis, we propose:

\begin{lstlisting}[language=Python,caption=Mario-DPO Training Configuration]
# DPO Training Configuration
model:
  base: gpt2-small (124M params)
  tokenizer: character-level (37 tokens)
  context_length: 2048

data:
    human_decisive_pairs: 816
    human_ties: 165 (optional weak preferences)
    synthetic_pairs: configurable, judge-filtered
    human_pair_weight: high
    synthetic_pair_weight: low

dpo:
  beta: 0.1          # KL penalty
  learning_rate: 1e-5
  epochs: 3
  batch_size: 32

inference:
  rejection_sampling_n: 10
  playability_filter: A* check
  style_conditioning: [STYLE: NINTENDO]
\end{lstlisting}

\subsection{Expected Outcomes}
\label{subsec:expected-outcomes}

If successful, Mario-DPO would:
\begin{enumerate}
    \item Generate levels achieving $>50\%$ win rate against Original in blind A/B tests
    \item Produce playable levels (A* verification) without post-hoc filtering
    \item Match the ``Nintendo style'' measured by low distance to the \texttt{original} corpus
    \item Enable style-conditioned generation via prompt tokens
\end{enumerate}

This work would represent a significant advancement in learned PCG, demonstrating that preference alignment can bridge the quality gap between procedural and hand-crafted content.

%==============================================================================
\section{Conclusions and Future Work}
\label{sec:conclusions}
%==============================================================================

\subsection{Summary of Contributions}
\label{subsec:contributions}

This project presents \textbf{PCG Arena}, a complete web platform for blind A/B testing of procedural content generators through human preference data. Key contributions include:

\begin{enumerate}
    \item \textbf{Browser-Based Mario Engine}: A faithful TypeScript port of the Mario AI Framework game engine, enabling gameplay without downloads.
    
    \item \textbf{Uncertainty-Aware Rating System}: Implementation of Glicko-2 providing not just rankings but confidence intervals essential for research validity.
    
    \item \textbf{Adaptive Matchmaking (AGIS)}: A novel algorithm optimizing battle selection for rapid rating convergence while ensuring coverage.
    
    \item \textbf{Research Infrastructure}: Comprehensive telemetry, data export, and analytics tools enabling research-grade data collection.
    
    \item \textbf{Open Platform}: Builder Profile system allowing researchers to submit and compete with their own generators.
    
    \item \textbf{Production Deployment}: Live platform at \url{https://pcg-arena.com} with proper security, authentication, and monitoring.
\end{enumerate}

\subsection{Limitations}
\label{subsec:limitations}

Several limitations should be acknowledged:

\begin{itemize}
    \item \textbf{Platform-Specific}: Currently limited to Mario-style platformer levels; generalizing to other game genres requires significant adaptation.
    
    \item \textbf{Desktop-Only}: No touch/mobile support limits the potential participant pool.
    
    \item \textbf{Selection Bias}: Self-selected participants may not represent general player preferences.
    
    \item \textbf{Level Length}: Fixed 30-second gameplay may not capture all relevant differences.
    
    \item \textbf{Single Metric}: Pairwise preference conflates multiple quality dimensions (fun, difficulty, aesthetics).
\end{itemize}

\subsection{Future Directions}
\label{subsec:future}

Several directions for future work emerge from this project:

\textbf{Short-term Improvements:}
\begin{itemize}
    \item Mobile/touch control support
    \item Tutorial and onboarding flow
    \item Social features (sharing results, challenging friends)
    \item Expanded tag vocabulary based on user feedback
\end{itemize}

\textbf{Research Extensions:}
\begin{itemize}
    \item Multi-armed bandit approaches for matchmaking
    \item Machine learning models predicting preference from level features
    \item Player clustering to identify preference types
    \item Cross-generator level morphing experiments
\end{itemize}

\textbf{Platform Generalization:}
\begin{itemize}
    \item Support for other game genres (puzzle games, roguelikes)
    \item Pluggable game engine architecture
    \item API for programmatic generator submission
    \item Federation with other PCG evaluation platforms
\end{itemize}

PCG Arena represents a step toward democratizing human evaluation in PCG research, providing the community with an accessible tool for empirically grounded generator comparison.

%==============================================================================
% References
%==============================================================================
\newpage
\bibliographystyle{plain}
% \bibliography{references}

\begin{thebibliography}{9}

\bibitem{glickman2012}
M.~E. Glickman, ``Example of the Glicko-2 system,'' 
\url{http://www.glicko.net/glicko/glicko2.pdf}, 2012.

\bibitem{marioai}
A.~Khalifa, M.~C.~Green, D.~Perez-Liebana, and J.~Togelius, ``General Video Game Level Generation,''
\emph{Proceedings of the Genetic and Evolutionary Computation Conference}, 2019.

\bibitem{summerville2018}
A.~Summerville, S.~Snodgrass, M.~Guzdial, C.~Holmg{\aa}rd, A.~K.~Hoover, A.~Isaksen, A.~Nealen, and J.~Togelius, ``Procedural Content Generation via Machine Learning (PCGML),''
\emph{IEEE Transactions on Games}, vol.~10, no.~3, pp.~257--270, 2018.

\bibitem{rafailov2023dpo}
R.~Rafailov, A.~Sharma, E.~Mitchell, S.~Ermon, C.~D. Manning, and C.~Finn, ``Direct Preference Optimization: Your Language Model is Secretly a Reward Model,''
\emph{Advances in Neural Information Processing Systems (NeurIPS)}, vol.~36, 2023.

\end{thebibliography}

\end{document}
